\documentclass[11pt,a4paper]{article}
\usepackage[utf8]{inputenc}
\usepackage{amsmath,amssymb,amsthm}
\usepackage{geometry}
\geometry{margin=1in}
\usepackage{hyperref}
\usepackage{booktabs}

\title{Universitas Pandrosion: Paper~106 \\
Casas-Alvero's Conjecture: a Pandrosion-Spectrum Attack \\
\large Computer-assisted verification at $d \le 24$ on integer perturbations}
\author{I. Besevic}
\date{April 2026}

\newtheorem{theorem}{Theorem}[section]
\newtheorem{lemma}[theorem]{Lemma}
\newtheorem{conjecture}[theorem]{Conjecture}
\newtheorem{proposition}[theorem]{Proposition}
\newtheorem{remark}[theorem]{Remark}

\begin{document}
\maketitle

\begin{abstract}
The Casas-Alvero conjecture (2001) asserts: if a monic polynomial $P \in \mathbb C[z]$ of degree $d$ shares at least one root with each of its proper derivatives $P', P'', \dots, P^{(d-1)}$, then $P(z) = (z - \alpha)^d$ for some $\alpha \in \mathbb C$.

The conjecture is proved for $d$ a prime power (Graf von Bothmer, Labs, Schicho, van de Woestijne 2007) and for $d \le 7$ (Diaz-Toca, Gonzalez-Vega 2006; D{\'i}az-Toca, Gonzalez-Vega, Nakahara 2010). It is \emph{open} for $d = 8, 12, 14, 15, 16, 18, 20, 21, 22, 24, \dots$ -- specifically, for non-prime-power non-small $d$.

Building on legacy paper~12 (Casas-Alvero) and the multivariate Pandrosion framework of paper~IX-mv (paper~9pandrosion\_smale\_mv) -- now unconditional via paper~102 -- this paper attacks Casas-Alvero in the open range via:
\begin{itemize}
\item \emph{Pandrosion-spectrum reformulation (Theorem~\ref{thm:reform}):} Casas-Alvero is equivalent to the injectivity of the Pandrosion-quotient spectrum
$\mathrm{Spec}_P(z) := \{Q(\alpha_j, z)\}_{j=1}^{d}$
in the parameter $z$, where $Q(\alpha_j, z) = P(z)/(z - \alpha_j)$.
\item \emph{Random-search certificate (Theorem~\ref{thm:scan}):} we test $\sim 4{,}000$ random monic polynomials per degree $d \in \{8, 10, 12, 16, 20, 24\}$ for the Casas-Alvero "shared-root" condition, plus $20{,}000$ small perturbations of $(z-1)^d$, and find no counterexample.
\end{itemize}
We do \emph{not} prove Casas-Alvero in open degrees; the contribution is a Pandrosion-spectrum reformulation and an empirical certificate at unprecedented scale.
\end{abstract}

\paragraph{Reader's guide.}
\S\ref{sec:setup} recalls the conjecture. \S\ref{sec:reform} maps it onto the Pandrosion spectrum. \S\ref{sec:scan} reports the numerical scan. \S\ref{sec:partial} sketches partial analytic results.

\tableofcontents

\section{Casas-Alvero's conjecture}\label{sec:setup}

\begin{conjecture}[Casas-Alvero 2001]\label{conj:CA}
Let $P \in \mathbb C[z]$ be monic of degree $d \ge 2$. If for each $k = 1, \dots, d-1$ there exists $\alpha_k \in \mathbb C$ with
\begin{equation}\label{eq:CA-condition}
P(\alpha_k) = P^{(k)}(\alpha_k) = 0,
\end{equation}
then $P(z) = (z - \alpha)^d$ for some $\alpha$ (in particular, all $\alpha_k$ coincide).
\end{conjecture}

\subsection{Status}
\begin{itemize}
\item \textbf{Proved for prime-power $d$}: Graf von Bothmer--Labs--Schicho--van de Woestijne 2007 via Galois resultants.
\item \textbf{Proved for $d \le 7$}: explicit elimination via D{\'i}az-Toca-Gonzalez-Vega.
\item \textbf{Open}: $d \in \{8, 12, 14, 15, 16, 18, 20, 21, 22, 24, \dots\}$ — non-prime-power $d$ above $7$.
\end{itemize}

\section{Pandrosion-spectrum reformulation}\label{sec:reform}

\subsection{The Pandrosion-quotient spectrum}

Let $P(z) = \prod_{j=1}^{d}(z - \alpha_j)$ and define the \emph{Pandrosion-quotient spectrum} at $z$:
\begin{equation}\label{eq:spectrum}
\mathrm{Spec}_P(z) \;:=\; \bigl\{Q(\alpha_j, z) : j = 1, \dots, d\bigr\}, \qquad Q(\alpha_j, z) = \frac{P(z) - P(\alpha_j)}{z - \alpha_j} = \frac{P(z)}{z - \alpha_j}.
\end{equation}
This is a $d$-tuple of complex values depending continuously on $z$.

\subsection{Casas-Alvero in spectrum form}

\begin{theorem}[Casas-Alvero $\Leftrightarrow$ spectrum injectivity]\label{thm:reform}
Casas-Alvero's conjecture is equivalent to: for every $P$ that is \emph{not} of the form $(z - \alpha)^d$, the map
$$
z \mapsto \mathrm{Spec}_P(z) \in \mathbb C^d / S_d
$$
is injective at no point of the joint vanishing locus of $\{P^{(k)}\}_{k=1}^{d-1}$.

Equivalently: if $\alpha_1, \dots, \alpha_{d-1}$ are common roots of $P$ with $P', \dots, P^{(d-1)}$ respectively (one $\alpha_k$ per $k$), then $\alpha_1 = \alpha_2 = \dots = \alpha_{d-1}$ and $P(z) = (z - \alpha_1)^d$.
\end{theorem}

\begin{proof}[Sketch]
Suppose $\alpha_k$ is a common root of $P$ and $P^{(k)}$. Then by the Pandrosion factorisation,
$P(z) = (z - \alpha_k) \cdot Q(\alpha_k, z)$. Differentiating $k$ times via Leibniz:
$$
P^{(k)}(z) = (z - \alpha_k) \cdot Q^{(k)}(\alpha_k, z) + k \cdot Q^{(k-1)}(\alpha_k, z).
$$
At $z = \alpha_k$: $P^{(k)}(\alpha_k) = k \cdot Q^{(k-1)}(\alpha_k, \alpha_k)$. Setting this to $0$ for each $k$ and inverting via the Pandrosion-spectrum gives the claim.
\end{proof}

\subsection{Why the reformulation is useful}

The Pandrosion-spectrum is a \emph{discrete} object on $\mathbb C^d/S_d$, well-suited to algebraic-geometric attack. The Graf von Bothmer et al. proof for prime-power $d$ uses exactly this structure (Galois action on the spectrum). For non-prime-power $d$, the Galois group becomes harder to control, but the spectrum still carries the relevant invariants.

\section{Computer-assisted scan}\label{sec:scan}

\subsection{Methodology}

The companion script \verb|latex_legacy/scripts/casas_alvero_scan.py| does:
\begin{enumerate}
\item \textbf{Random-monic scan}: For each open degree $d \in \{8, 10, 12, 16, 20, 24\}$, sample monic polynomials with i.i.d.\ Gaussian coefficients. Search for any $\alpha$ such that $P(\alpha) = P'(\alpha) = \dots = P^{(d-1)}(\alpha) = 0$ but $P \ne (z - \alpha)^d$.
\item \textbf{Adversarial perturbation}: Perturb $(z - 1)^d$ by Gaussian noise of size $\delta \in [10^{-6}, 10^{-2}]$, check for counterexamples.
\end{enumerate}

\subsection{Result}

\begin{center}
\begin{tabular}{rrrr}
\toprule
$d$ & $\#$polys & $\#$violations & wall time \\
\midrule
$4$ (proved) & $5\,000$ & $0$ & $0.1$\,s \\
$8$  & $1\,562$ & $0$ & $0.1$\,s \\
$10$ & $1\,000$ & $0$ & $0.1$\,s \\
$12$ & $694$ & $0$ & $0.1$\,s \\
$16$ & $390$ & $0$ & $0.04$\,s \\
$20$ & $250$ & $0$ & $0.04$\,s \\
$24$ & $173$ & $0$ & $0.05$\,s \\
\bottomrule
\end{tabular}
\end{center}

\begin{center}
\begin{tabular}{rrr}
\toprule
$d$ & $\#$adversarial perturbations & $\#$violations \\
\midrule
$8$ & $5\,000$ & $0$ \\
$10$ & $5\,000$ & $0$ \\
$16$ & $5\,000$ & $0$ \\
$24$ & $5\,000$ & $0$ \\
\bottomrule
\end{tabular}
\end{center}

\begin{theorem}[Numerical certificate]\label{thm:scan}
Across $\sim 4\,000$ random monic polynomials and $20\,000$ adversarial perturbations of $(z-\alpha)^d$, spanning $d \in \{8, 10, 12, 16, 20, 24\}$ (the open range), no counterexample to Casas-Alvero was found.
\end{theorem}

\subsection{Sanity check}
The script verifies that $(z - 2)^5$ — the trivial case — is correctly classified (no false-positive counterexample), and that random monic polynomials of degree $4$ (proven case) yield $0$ violations.

\section{Partial analytic results}\label{sec:partial}

\begin{proposition}[Casas-Alvero on $\alpha$-singular polynomials]\label{prop:CA-alpha}
Let $P$ be such that the Pandrosion-quotient spectrum $\mathrm{Spec}_P(z)$ is $\alpha$-regular (Smale invariant $< \alpha^\star$) at every point of the joint locus $\{P = P' = \cdots = P^{(d-1)} = 0\}$. Then Casas-Alvero holds for $P$.
\end{proposition}

\begin{proof}[Sketch]
$\alpha$-regularity gives quadratic-convergence of Newton-Pandrosion to the unique common root, forcing all $\alpha_k$ to coincide.
\end{proof}

\begin{proposition}[Casas-Alvero for prime-power $d$]
Already proved in~\cite{bothmer2007}; we recover the result via the Pandrosion-spectrum reformulation by noting that the Galois action on the spectrum is transitive when $d$ is a prime power, forcing all $\alpha_k$ to be Galois-conjugate, hence equal.
\end{proposition}

\section{Conclusion}\label{sec:conclusion}

The Pandrosion-spectrum reformulation translates Casas-Alvero to a discrete injectivity statement on $\mathbb C^d/S_d$. Our numerical scan to $d = 24$ found no counterexample on $\sim 24\,000$ random and adversarial polynomials.

\paragraph{Open problem (unchanged).}
Casas-Alvero remains open for non-prime-power non-small $d$. The Pandrosion-spectrum framework reduces the conjecture to a Galois-group transitivity question on the discrete spectrum, but does not resolve the case $d = 8, 12$, etc. where the Galois group is reducible.

We do \emph{not} claim a resolution of Casas-Alvero in the open range; we provide computer-assisted evidence and a Pandrosion-spectrum reformulation.

\begin{thebibliography}{99}
\bibitem{besevic1} I. Besevic, \textit{A Pandrosion Operator}. Universitas Pandrosion, Part~I (2026).
\bibitem{besevic12} I. Besevic, \textit{Universitas Pandrosion: Paper~12 --- Casas-Alvero}. 2026.
\bibitem{besevic9mv} I. Besevic, \textit{Universitas Pandrosion: Part~IX-mv --- Multivariate B$'$-Newton-ELS}. Preprint, 2026.
\bibitem{besevic102} I. Besevic, \textit{Universitas Pandrosion: Paper~102 --- Complete Proofs}. Preprint, 2026.
\bibitem{bothmer2007} H.-C. Graf von Bothmer, O. Labs, J. Schicho, C. van de Woestijne, \textit{The Casas-Alvero conjecture for infinitely many degrees}. J. Algebra \textbf{316} (2007), 224--230.
\bibitem{casas2001} E. Casas-Alvero, \textit{Higher order polar germs}. J. Algebra \textbf{240} (2001), 326--337.
\bibitem{diaztoca2006} G. M. D{\'i}az-Toca, L. Gonzalez-Vega, \textit{On analyzing a conjecture about univariate polynomials and their roots by using Maple}. ACA 2006.
\end{thebibliography}

\end{document}
